\chapter*{Table of Glossary}
\addcontentsline{toc}{chapter}{Table of Glossary}

\begin{description}

\item[A1 (Linked brushing)] An invariant of the dashboard: within
a single card, selecting an entity in one chart highlights it in
every other chart of the same card. Implementation uses
Streamlit \texttt{session\_state} and \texttt{st.rerun()}.

\item[A2 (Story subtitle)] An invariant of the dashboard: every
chart carries a descriptive title \emph{and} a one-sentence story
subtitle that tells the viewer what to notice in the chart.

\item[A3 (Universal Entity Detail Panel)] An invariant of the
dashboard: every clickable entity (route, stop, school, company,
day, hour, hex cell, vehicle, driver) opens a side panel of
approximately forty percent of viewport with a multi-chart
mini-dashboard about that entity.

\item[AFC] Automated Fare Collection. The HCMC ticketing system
that logs each boarding event, the fare class, and the payment
channel. The local AFC dataset spans 1~January to 31~March~2026
with 17.8 million boarding records.

\item[Aggregate view] A pre-computed \texttt{CREATE OR REPLACE TABLE}
in DuckDB that serves one or more dashboard charts. The
warehouse defines eight aggregate views; every chart maps to
exactly one of them (the chart-to-aggregate data contract).

\item[BusLens] The product name of the visual analytics platform
delivered in this thesis.

\item[Card] One of the six themes in the dashboard hub
(School-Aware Demand, Demand \& Crowding, Spatial \& Connectivity,
Financial \& Operations, Network Health \& Performance, Routes).
Each card is anchored internally to one or two user roles but
presents the same chrome to every user.

\item[deck.gl HexagonLayer] A spatial-aggregation primitive that
bins boardings into hex cells of configurable radius
(200\,m / 500\,m / 1\,km) and colour-codes them by density.

\item[\texttt{dim\_route} canonical key] The AFC \texttt{route\_no}
field is the canonical \texttt{route\_id}; the GPS \texttt{route\_no}
is stored as a join field on \texttt{dim\_route}. Routes are
flagged as GPS-only or AFC-only where they appear in only one
source.

\item[Ghost run] A trip recorded in schedule data where GPS
positioning shows door openings on fewer than 20\% of scheduled
stops. Used as a proxy for operational waste in
Chart 4.1.

\item[Ghost stop] A stop that appears in the schedule but where
GPS logs show door openings on fewer than 20\% of scheduled
visits. Used in Chart 3.4.

\item[Gridify] A decoupling of data and visual grid abstractions
from Liu~(2022), used in the dashboard's hex-grid map layer.

\item[GroupSet / UpSet matrix] A set-intersection visualisation
from Liu~(2022) used in Chart~6.2 to show routes that belong to
multiple archetypes simultaneously (e.g., ``school AND
commuter'').

\item[HCMC Bus Analytics] The brand string shown in the dashboard
top bar.

\item[Hub $\rightarrow$ Card $\rightarrow$ Chart] The three-level
information-architecture hierarchy of the dashboard.

\item[Marey time--distance diagram] Chart 6.5; each trip on a
single route plotted as a diagonal line on a time-by-stop-distance
grid. Source: Lin~\&~Zeng (\cite{linzeng2014}).

\item[MCV (Multiple Coordinated Views)] An interaction model from
Liu~(2022) that the A1 linked-brushing invariant implements.

\item[QoS-1 / QoS-2 / QoS-3] Three Quality-of-Service tiers from
Srisakaokul~\textit{et~al.}~(\cite{srisakaokul2023}): QoS-1 = trip
completion, QoS-2 = on-path driving (50\,m route buffer), QoS-3
= on-schedule within tolerance~$T$.

\item[SDVF tiers] Solution-Decision-Value-Framework prioritisation
applied across the design space: FLAGSHIP (the differentiator) /
MUST (table-stakes) / SHOULD (depth) / COULD (secondary) / WON'T
(yet, scoped out for Phase~1).

\item[Stories layer] A scrollytelling section of the hub
presenting four curated narratives (The School Story, The
T\d ết~Story, The Money Story, The Network Map) on the same chart
canvases the cards expose.

\item[Threshold $T$] A tunable minutes-of-delay parameter for
anomaly detection in Chart 5.5. Pre-computed at $T \in
\{1, 2, 5, 10, 15\}$ minutes so the slider snaps between levels.

\end{description}
